%%======================================================================
\section{Introduction}
\label{sec:intro}
%%======================================================================

Full-size bipedal humanoids can now execute clinical procedures in the
laboratory. A Unitree G1 teleoperated by a non-clinician completed seven
distinct procedures: auscultation, ultrasound, intravenous catheter insertion,
and four others~\cite{atar2025humanoids}; a second system performed
laparoscopic surgery on a cadaveric torso~\cite{liang2025lapsurgie}; a third
acted as a surgical first assistant in a cadaveric bilateral
sphenoidectomy~\cite{cho2026humanoid}. Not one of these systems holds U.S. Food and Drug Administration (FDA)
clearance or Conformit\'{e} Europ\'{e}enne (CE) marking for any clinical
application. No registered clinical
trial has enrolled a patient for a study involving a full-size bipedal humanoid
platform. We nonetheless expect this gap between laboratory capability and
clinical evidence to narrow, because two developments converged after
2022~\cite{gu2025humanoid,cao2025humanoid}. On the hardware side, at least
thirteen commercial full-size platforms were released between 2022 and 2026
(Section~\ref{sec:systems}), reaching the payload, dexterity, and locomotion
thresholds needed to attempt clinical tasks in human-scale
environments~\cite{ifr2025humanoid,unitree2024g1,figureai2024unveils}. On the
software side, the maturation of electric actuation and learning-based
locomotion was accompanied by large language models (LLMs), vision-language
models (VLMs), and vision-language-action (VLA) models that reduced the
task-specific programming which had previously limited general-purpose humanoid
deployment~\cite{firoozi2025foundation,kawaharazuka2024realworld,kim2024survey}.
Where earlier systems had to be re-engineered for each new task, instruction
following, task generalization, and adaptation to unstructured environments
became feasible with far less bespoke
engineering~\cite{argall2009survey,ravichandar2020recent}. In roughly three
years, the field moved from proof-of-concept locomotion to laboratory execution
of clinical procedure
sequences~\cite{radosavovic2024realworld,atar2025humanoids,cho2026humanoid}.

Several recent reviews approach this subject from adjacent vantage points, and we
position our contribution against them along three axes: scope, method, and
output. In \emph{scope}, Cao~\cite{cao2025humanoid} reviews humanoid robots and
humanoid AI as a general technology, and Liu et al.~\cite{liu2025screens} survey
embodied AI across the full breadth of healthcare, from diagnostic agents to
virtual assistants; neither treats the full-size bipedal humanoid as a distinct
clinical device class defined by its form factor. In \emph{method}, both are
broad narrative syntheses, whereas we run a systematic search with explicit
inclusion and exclusion criteria (detailed below) that isolates every study
placing such a platform in a clinical task. In \emph{output}, we pair
that evidence with a platform-by-platform capability assessment, a regulatory
pathway analysis, and per-domain technology-readiness scoring, none of which the
prior surveys provide. Cunha et al.~\cite{cunha2025humanoid} is the closest
dedicated work, a 2025 scoping review, but its November 2023 search cutoff
predates the LLM/VLA integration wave. Earlier healthcare-robot
reviews~\cite{azeta2018review,mukherjee2022humanoid,ozturkcan2022humanoid,huang2023intelligent,morgan2022robots}
treat platform types interchangeably, which obscures the form-factor-specific
barriers that in practice determine clinical deployability. The 2022--2026
bipedal-humanoid healthcare literature therefore still lacks a dedicated
systematic review.

This review covers bipedal, free-standing humanoid robots of approximately
150--200~cm height designed for general-purpose manipulation and locomotion in
human-scale environments, with literature from January 2022 to March 2026.
Wheeled social robots (Pepper, ARI), sub-100~cm platforms (NAO), surgical arms
(da~Vinci), and powered exoskeletons are excluded; each exclusion is argued on
form-factor grounds in Section~\ref{sec:background}. RHP
Friends~\cite{benallegue2025rhp} is the only non-commercial academic platform
retained without caveat. The systematic search ran across five databases: arXiv
(cs.RO, cs.AI, eess.SY), PubMed, IEEE Xplore, ACM Digital Library, and Google
Scholar. The primary query combined platform and application terms:
(\emph{humanoid robot} OR \emph{bipedal robot}) AND (\emph{healthcare} OR
\emph{medical} OR \emph{clinical} OR \emph{nursing} OR \emph{rehabilitation} OR
\emph{surgery}), supplemented by domain-specific secondary searches and
citation-list hand-search. Title and abstract screening of 147 candidate papers
yielded 52 full-text reviews; applying inclusion and exclusion criteria produced
a focused set of 14 clinical studies. To our knowledge this is the complete
clinical evidence base for bipedal humanoid research to date, and mapping it in
full, rather than sampling it, is itself one contribution of this review; the
larger body of material we survey spans platforms, software, safety standards,
and regulation.
Concurrent with the literature search, we surveyed 13
commercial and 12 academic systems meeting scope criteria, verifying each
specification against manufacturer documentation or peer-reviewed publications;
the full list, per-platform sources, and specifications are given in
Table~\ref{tab:platforms}~\cite{figureai2024unveils,unitree2024g1,bostondynamics2024atlas,apptronik2023apollo,sanctuaryai2023phoenix,dafarra2024icub3}. This review makes four contributions:


\begin{itemize}
\item A systematic review of humanoid robots in healthcare across the 2022--2026
  period, over which electric actuation, learning-based control, and LLM/VLA
  integration matured in parallel, with Technology Readiness Level (TRL)
  assignments per ISO~16290:2013.
\item A capability assessment of 25 bipedal humanoid platforms (13 commercial,
  12 academic) against five clinical domain requirements.
\item A regulatory pathway analysis covering FDA De~Novo, EU Medical Device
  Regulation (MDR) 2017/745, and EU AI~Act 2024/1689.
\item A Humanoid Clinical Deployment Readiness (HCDR) framework linking technical
  maturity to regulatory pathway for five clinical roles.
\end{itemize}

The remainder of this paper is organized as follows:
Section~\ref{sec:background} provides background; Section~\ref{sec:systems}
surveys current humanoid capabilities; Section~\ref{sec:applications} analyzes
what healthcare demands from humanoid robots; Section~\ref{sec:challenges}
examines clinical deployment challenges; Section~\ref{sec:discussion} discusses
future directions; Section~\ref{sec:conclusion} concludes.

